\uuid{TNdQ}
\exo7id{4239}
\titre{exo7 4239}
\auteur{quercia}
\organisation{exo7}
\datecreate{2010-03-12}
\isIndication{false}
\isCorrection{true}
\chapitre{Calcul d'intégrales}
\sousChapitre{Autre}
\cpge{mpsi-mp2i.integration.c, pcsi.integration.b, ptsi.integration.b, mpsi-mp2i.nombres-reels-suites.d, pcsi.nombres-reels-suites.c, ptsi.nombres-reels-suites.c, bcpst1.derivees-primitives-integrales.d, bcpst1.integration-segment.a, mp-mpi.integration-intervalle.b, pc.integration.a, psi.integration.a, bcpst2.integrales-generalisees}
\module{Analyse}
\niveau{L1}
\difficulte{}

\contenu{
\question{
Calculer pour tout $n \in \N^*$~:
$I_n =  \int_{x=0}^\pi \frac{dx}{1+\cos^2(nx)}$.
}
\reponse{
Couper en intervalles de $k\pi/n$.
         On obtient $I_n = \frac\pi{\sqrt2}$ pour tout $n\ge 1$.
}
}
